
We measured GRPO advantage variance in two conditions—function-level execution feedback on competitive programming tasks and repository-level execution feedback on software engineering tasks—using CodeLlama-7b-Instruct-hf with identical GRPO configuration (G=8, B=4, 120 steps). The result is a 76× advantage variance gap (0.004 vs. 0.317, p=5.34×10⁻⁴⁴, Cohen's d=1.904), explained by the difference in positive reward rates (≈0\% vs. ≈6\%). When positive reward rate is near zero, GRPO groups are degenerate and gradient contributions are zero.

We additionally implemented and smoke-tested training infrastructure for a 4-condition curriculum GRPO experiment on APPS+CodeContests. All four conditions (curriculum, uniform, easy_only, hard_only) completed the 10-step smoke test with exit code 0. Reward density was 0.0 in all conditions throughout the smoke test, consistent with the reward collapse finding. Full training was not executed; no behavioral comparison of curriculum vs. uniform sampling is available.

The primary open question raised by this work is whether assumption A1—that easy-tier APPS problems are solvable at initialization for DeepSeek-Coder-7B-base—holds. If it does, the full curriculum GRPO experiment is warranted. If it does not, alternative interventions such as SFT warm-up or partial credit reward shaping are needed before curriculum ordering can be expected to produce a gradient signal advantage.

Three concrete directions follow from these findings:

\begin{enumerate}
\item Measure per-difficulty-tier solve rates at initialization for the target model (DeepSeek-Coder-7B-base on APPS tiers 0–4) before running full curriculum GRPO training.
\item Implement and compare partial credit reward functions for function-level tasks to test whether reward engineering is a more tractable intervention than curriculum ordering for the advantage collapse problem.
\item Design a within-function-level experiment comparing easy-tier vs. hard-tier APPS advantage variance to directly test the curriculum mechanism rather than the cross-granularity confound.
\end{enumerate}
